\documentclass[11pt,oneside]{article}
\usepackage{geometry}
\geometry{letterpaper}
\usepackage{amssymb}
\usepackage{amsmath}
\usepackage{mathpazo,euler}
\usepackage{hyperref}

\title{Automating Sylow's Theorems in Haskell}
\date{}

\begin{document}
\maketitle

\section{Overview}
This document describes a Haskell implementation of a specialized
theorem prover for Sylow-style arguments in finite group theory. The
program automates common reasoning patterns that show no simple group
of a given order can exist. The implementation encodes those techniques
as rewrite rules and searches for a contradiction, while using static
types and a clear module structure.

\section{Module Map}
The code is split into focused modules:
\begin{center}
\begin{tabular}{ll}
\texttt{src/Core.hs} & Core data types: \texttt{Fact}, \texttt{Arg}, labels, theorems \\
\texttt{src/Predicates.hs} & Fact constructors for group-theoretic predicates \\
\texttt{src/Env.hs} & Proof environment and pure state transitions \\
\texttt{src/EnvPrint.hs} & Printing of proof traces (IO only) \\
\texttt{src/Auto.hs} & Proof search loop and template matching \\
\texttt{src/Theorems.hs} & Sylow techniques as theorems and rules \\
\texttt{src/NumberTheory.hs} & Primes, divisors, and factorization helpers \\
\texttt{src/DisjunctionHash.hs} & Stable disjunction labeling via hashing \\
\texttt{app/Main.hs} & CLI entry point (interactive and batch) \\
\end{tabular}
\end{center}

\section{Facts, Arguments, and Matching}
Facts are predicate names with ordered arguments. Arguments are typed
to make matching explicit:
\begin{itemize}
\item \texttt{Sym s}: a concrete symbol \texttt{s} (treated literally)
\item \texttt{Var v}: a template variable \texttt{v} (can match anything)
\item \texttt{Exact e}: a literal that must match \texttt{e} (e.g. \texttt{1})
\item \texttt{Fresh f}: a new symbol to be generated by the environment
\end{itemize}

Template matching unifies \texttt{Var} and \texttt{Fresh} with concrete
symbols and requires \texttt{Exact} to match literally. After a theorem
fires, any \texttt{Fresh} arguments are replaced with globally fresh
symbols so that new objects do not collide with existing names.

\section{Disjunction Handling}
Disjunctions represent logical \textit{or} at the fact level. When a
disjunction is produced, each branch is added as a fact, and each such
fact records which branch it depends on. A goal is considered proven
only if every full assignment of disjunction branches is covered by a
goal derivation. This is implemented by tracking disjunction ancestry
sets and checking coverage over the Cartesian product of branches.

\section{Theorem Set}
The technique library mirrors standard Sylow-based reasoning:
\begin{itemize}
\item Sylow subgroup counts and existence.
\item Normality from a unique Sylow subgroup.
\item Embedding into alternating groups and index arguments.
\item Counting arguments for elements of $p$-power order.
\item Normalizer-of-intersection arguments and related constraints.
\end{itemize}
Each technique is a theorem object that matches a template and yields
new facts or disjunctions, allowing the proof engine to chain them
automatically.

\section{Worked Example (Order 12)}
Assume a simple group $G$ of order $12$. The prover:
\begin{enumerate}
\item Applies Sylow to derive possible counts of Sylow $2$- and
  $3$-subgroups.
\item If a Sylow count is $1$, it derives that the subgroup is normal
  and hence $G$ is not simple.
\item In remaining branches, it uses counting arguments to show the
  number of elements of $2$- and $3$-power order exceeds $12$.
\end{enumerate}
In each branch, a contradiction is derived, so the goal \texttt{false}
is proven.

\section{Limitations and Scope}
This system is intentionally specialized. It is not a full first-order
theorem prover, and it does not encode group axioms. The search is
brute-force over the technique library, so disjunctions can cause
combinatorial blowups for large orders. Nevertheless, for many orders
the technique set is sufficient and produces short, readable proofs.

\section{References}
Useful references for the underlying group theory include standard
texts such as:
\begin{itemize}
\item Dummit and Foote, \textit{Abstract Algebra}
\item Isaacs, \textit{Finite Group Theory}
\end{itemize}

\end{document}
